% Kompilieren: latexmk -pdf noposroots.tex
% Alternativ: pdflatex noposroots; biber noposroots; pdflatex noposroots; pdflatex noposroots
% Literatur: biblatex mit Biber. Nur zitierte Eintraege werden ausgegeben.

\documentclass[a4paper,11pt]{article}

% --- Schrift und Seitenlayout ---
\usepackage[T1]{fontenc}
\usepackage{lmodern} % Skalierbare Schrift, nahe an der AMS-Standardschrift.
\usepackage[english]{babel} % Fuer deutsche Notizen: english durch ngerman ersetzen.
\usepackage{microtype}
\usepackage[margin=2.5cm]{geometry}

% --- Schlichtes Notizlayout ---
\pagestyle{plain} % Nur Seitenzahlen; keine laufenden Kopfzeilen.
\setlength{\parindent}{0pt}
\setlength{\parskip}{0.35em plus 0.1em minus 0.1em}


% --- Literatur mit biblatex und Biber ---
\usepackage{csquotes}
\usepackage[
backend=biber,
style=numeric-comp,
sorting=nyt,
maxbibnames=99
]{biblatex}
\addbibresource{literature.bib}
\renewcommand*{\bibfont}{\small}


% --- Mathematik und Listen ---
\usepackage{mathtools,amssymb,amsthm}
\usepackage{enumitem}
\newlist{compactitem}{itemize}{3}
\setlist[compactitem]{topsep=2pt,partopsep=0pt,itemsep=0pt,parsep=0pt}
\setlist[compactitem,1]{label=\textbullet}
\setlist[compactitem,2]{label=--}
\setlist[compactitem,3]{label=*}

% Bei Bedarf aktivieren:
% \usepackage{graphicx}               % Bilder mit \includegraphics
% \usepackage{tikz-cd}                % Kommutative Diagramme
% % \usepackage[notref,notcite]{showkeys} % Fuer die Fehlersuche aktivieren. Labelnamen am Rand
% \usepackage{newpxtext,newpxmath}     % Palatino-Schrift wie in notes2

% --- Links und Querverweise ---
\usepackage{xcolor}
\definecolor{darkred}{RGB}{160,0,0}
\definecolor{darkblue}{RGB}{0,0,160}
\definecolor{lightpurple}{RGB}{140,50,150}
\usepackage[colorlinks=true,linkcolor=darkblue,citecolor=lightpurple,
urlcolor=darkblue]{hyperref}
\usepackage[sort,nameinlink]{cleveref}
\usepackage{thmtools}

% \usepackage{extarrows} % Bei Bedarf fuer zusaetzliche lange Pfeile.

% --- Satzumgebungen ---
% Gemeinsame fortlaufende Nummerierung: Theorem 1, Lemma 2, Example 3, ...
% Aufrechte Schrift wie in sbei-no-ai; fuer kursive Saetze style=plain setzen.
% thmtools erhaelt bei \cref und \Cref die richtigen Umgebungsnamen.
\declaretheorem[style=definition,name=Theorem]{theorem}
\declaretheorem[style=definition,name=Lemma,sibling=theorem]{lemma}
\declaretheorem[style=definition,name=Proposition,sibling=theorem]{proposition}
\declaretheorem[style=definition,name=Corollary,sibling=theorem]{corollary}
\declaretheorem[style=definition,name=Definition,sibling=theorem]{definition}
\declaretheorem[style=definition,name=Example,sibling=theorem]{example}
\declaretheorem[style=definition,name=Remark,sibling=theorem]{remark}
\declaretheorem[style=definition,name=Conjecture,sibling=theorem]{conjecture}
\declaretheorem[style=definition,name=Question,sibling=theorem]{question}
\declaretheorem[style=definition,name=Notation,sibling=theorem]{notation}
\declaretheorem[style=definition,name=Claim,sibling=theorem]{claim}
\declaretheorem[style=definition,name=Note,sibling=theorem]{note}
\declaretheorem[style=definition,name=Convention,sibling=theorem]{convention}
% Gleichungen werden ebenfalls fortlaufend nummeriert: (1), (2), ...
% Fuer deutsche Notizen ausserdem hier name=... anpassen (z.B. name=Satz).

% --- Haeufige Befehle ---
\newcommand{\NN}{\mathbb{N}}
\newcommand{\ZZ}{\mathbb{Z}}
\newcommand{\QQ}{\mathbb{Q}}
\newcommand{\RR}{\mathbb{R}}
\newcommand{\CC}{\mathbb{C}}
\newcommand{\PP}{\mathbb{P}}
\newcommand{\FF}{\mathbb{F}}
\newcommand{\fr}[1]{\mathfrak{#1}}
\newcommand{\mideal}{\mathfrak{m}}
\newcommand{\nideal}{\mathfrak{n}}
\newcommand{\ol}[1]{\overline{#1}}
\newcommand{\card}[1]{\lvert#1\rvert}
\newcommand{\Set}[1]{\left\{\,#1\,\right\}}
\newcommand{\defas}{\coloneqq}
\newcommand{\asdef}{\eqqcolon}
\newcommand{\isoto}{\xrightarrow{\sim}}

\DeclareMathOperator{\im}{im}
\DeclareMathOperator{\coker}{coker}
\DeclareMathOperator{\id}{id}
\DeclareMathOperator{\rank}{rk}
\DeclareMathOperator{\Hom}{Hom}
\DeclareMathOperator{\tor}{Tor}
\DeclareMathOperator{\Ext}{Ext}
\DeclareMathOperator{\cone}{cone}
\DeclareMathOperator{\Spec}{Spec}
\DeclareMathOperator{\reg}{reg}
\DeclareMathOperator{\pd}{pd}
\DeclareMathOperator{\depth}{depth}
\DeclareMathOperator{\height}{height}
\DeclareMathOperator{\ch}{char}
\DeclareMathOperator{\lcm}{lcm}
\DeclareMathOperator{\rad}{rad}
\DeclareMathOperator{\initial}{in}
\DeclareMathOperator{\hilb}{Hilb}
\DeclareMathOperator{\supp}{supp}
\DeclareMathOperator{\zd}{ZD}

% Offene Stellen: \todo{Hier fehlt noch ein Argument.}
\newcommand{\todo}[1]{{\sffamily\color{darkred}\textbf{TODO:} #1}}

% \usepackage[notref,notcite]{showkeys} % Fuer die Fehlersuche aktivieren.

% --- Titel ---
\title{No positive roots} % Hier das Thema einsetzen.
\author{Niclas Niederdrenk}
\date{October 15, 2026} % Festes Datum; fuer das aktuelle Datum: \date{\today}.

% Kompakter, linksbuendiger Titel mit Autor und Datum.
\makeatletter
%\renewcommand{\maketitle}{%
%	\begingroup
%	\setlength{\parskip}{0pt}%
%	{\LARGE\bfseries \@title\par}%
%	\vspace{0.45em}%
%	{\small\color{black!65}\@author\enspace\textperiodcentered\enspace\@date\par}%
%	\endgroup
%	\vspace{1.1em}%
%}
\makeatother

\begin{document}
	\maketitle
	% \tableofcontents % Bei laengeren Notizen aktivieren.
	
	% Ab hier eigener Inhalt. Das kurze Beispiel zeigt die wichtigsten Befehle.
	
	
	The following theorem appears as Exercise~3(a) in Chapter~3 of \cite{Stanley2018}, which covers random walks. I solved the exercise while reading the book in August 2025, though it took me quite a while. These notes present my argument, with the aim of using it as a test case for AI-assisted formalization and verification in Lean.
	
	\begin{theorem}\label{main-theorem}
		Let $p \in \RR[x]$ be a nonzero polynomial. The following conditions
		are equivalent:
		\begin{enumerate}
			\item There exists a nonzero polynomial $q \in \RR[x]$ such that
			all coefficients of $pq$ are nonnegative.
			\item The polynomial $p$ has no positive real root, i.e.,
			$p(\alpha) \neq 0$ for every $\alpha > 0$.
		\end{enumerate}
	\end{theorem}
	
	The implication "$\Rightarrow$" is immediate: a nonzero polynomial with
	nonnegative coefficients has no positive real root, and any root of $p$
	is also a root of $pq$. It therefore remains to prove the implication
	"$\Leftarrow$".
	
	If $p$ is constant, we may simply take $q=1/p$. Otherwise, we can write
	$
	p=c p_1 \cdots p_n,
	$
	where $c\in\RR\setminus\{0\}$ and each $p_i$ is a monic irreducible
	polynomial of degree one or two. If $p$ has no positive real root, then
	neither does any of its irreducible factors.
	
	The product of polynomials with nonnegative coefficients again has
	nonnegative coefficients. Thus, if for each $i$ we can find a nonzero
	polynomial $q_i$ such that $p_iq_i$ has nonnegative coefficients, then
	$
	q=c^{-1}q_1 \cdots q_n
	$
	satisfies
	$
	pq=p_1q_1 \cdots p_nq_n,
	$
	which has nonnegative coefficients. It therefore suffices to prove the
	claim for monic irreducible polynomials of degree one or two.
	
	The linear case is immediate. Let $p=x+a$. Its unique root is $-a$,
	so $p$ has no positive real root if and only if $a\geq0$.
	Thus, all coefficients of $p$ are already nonnegative, and we may
	take $q=1$.
	
	For the quadratic case, consider a monic irreducible polynomial
	$p=x^2+bx+a\in\RR[x]$. Since $p$ is irreducible over $\RR$, its
	discriminant satisfies
	$
	\Delta=b^2-4a<0.
	$
	In particular, $0\leq b^2<4a$, so $a>0$. If $b\geq0$, then $p$
	already has nonnegative coefficients, and we may take $q=1$.
	It therefore remains to prove the following lemma.
	
	\begin{lemma}\label{quadratic-case}
		Let $p=x^2+bx+a\in\RR[x]$ satisfy $a>0$, $b<0$, and
		$b^2-4a<0$. Then there exists a nonzero polynomial
		$q\in\RR[x]$ such that all coefficients of $pq$ are nonnegative.
	\end{lemma}
	
	We define a real sequence $(y_k)_{k\geq0}$ by setting $y_0=0$ and
	requiring the formal power series identity
	\[
	\frac{a}{x^2+bx+a}
	= \sum_{k=1}^{\infty} y_k x^{k-1}.
	\]
	Since $a\neq0$, this uniquely determines the sequence.
	Comparing coefficients, we obtain the linear recurrence
	\[
	y_0=0,\qquad
	y_1=1,\qquad
	y_{k+2}=-\frac ba\,y_{k+1}-\frac1a\,y_k
	\quad(k\geq0).
	\]
	
	For each $n\geq1$, define the polynomial
	$$
	q_n(x)=\sum_{k=1}^{n}y_kx^{k-1}.
	$$
	
	\begin{lemma}\label{pq-product}
		For every $n\geq1$, we have
		$
		(a+bx+x^2)q_n(x)
		=a-ay_{n+1}x^n+y_nx^{n+1}.
		$
	\end{lemma}
	
	\begin{proof}
		We argue by induction on $n$. For $n=1$, we have $q_1=1$,
		$y_1=1$, and $y_2=-b/a$. Hence
		\[
		(a+bx+x^2)q_1(x)
		=a+bx+x^2
		=a-ay_2x+y_1x^2.
		\]
		
		Now suppose that the identity holds for some $n\geq1$.
		Since
		\[
		q_{n+1}(x)=q_n(x)+y_{n+1}x^n,
		\]
		the induction hypothesis gives
		\begin{align*}
			(a+bx+x^2)q_{n+1}(x)
			&=a-ay_{n+1}x^n+y_nx^{n+1}
			+(a+bx+x^2)y_{n+1}x^n\\
			&=a+(y_n+by_{n+1})x^{n+1}
			+y_{n+1}x^{n+2}.
		\end{align*}
		The recurrence relation implies
		\[
		y_n+by_{n+1}=-ay_{n+2}.
		\]
		Consequently,
		\[
		(a+bx+x^2)q_{n+1}(x)
		=a-ay_{n+2}x^{n+1}+y_{n+1}x^{n+2},
		\]
		which completes the induction.
	\end{proof}
	
	To prove \Cref{quadratic-case}, it therefore suffices to show that
	$y_k\leq0$ for some $k\geq1$. Indeed, choose $n$ such that $y_{n+1}$
	is the first nonpositive term after $y_0$. Since $y_1=1$, we have
	$n\geq1$ and $y_n>0$. The identity above then shows that all
	coefficients of $pq_n$ are nonnegative. Moreover, $q_n\neq0$,
	since its constant coefficient is $y_1=1$.
	
	It remains to establish the following lemma.
	
	\begin{lemma}\label{nonpositive-member}
		There exists an integer $k\geq1$ such that $y_k\leq0$.
	\end{lemma}
	
	
	\begin{proof}
		The characteristic polynomial of the recurrence is
		\[
		\chi(t)=t^2+\frac ba\,t+\frac1a.
		\]
		Its discriminant satisfies
		\[
		\Delta'=\frac{b^2-4a}{a^2}<0,
		\]
		so its roots are distinct and complex conjugates:
		\[
		\lambda_{1,2}
		=\frac{-b\pm i\sqrt{4a-b^2}}{2a}.
		\]
		Since $a>0$ and $b<0$, we have
		\[
		\operatorname{Re}(\lambda_1)=\frac{-b}{2a}>0,
		\qquad
		\operatorname{Im}(\lambda_1)
		=\frac{\sqrt{4a-b^2}}{2a}>0.
		\]
		Therefore, we may write
		\[
		\lambda_1=re^{i\phi},
		\qquad
		\lambda_2=re^{-i\phi}
		\]
		for some $r>0$ and $0<\phi<\pi/2$.
		
		The general solution of the recurrence has the form
		\[
		y_k=c\lambda_1^k+d\lambda_2^k
		\qquad(k\geq0)
		\]
		for some $c,d\in\CC$. The initial conditions $y_0=0$ and $y_1=1$
		give
		\[
		c+d=0,\qquad c\lambda_1+d\lambda_2=1.
		\]
		Consequently, $d=-c$ and
		\[
		c=\frac{1}{\lambda_1-\lambda_2}
		=\frac{1}{2i\,\operatorname{Im}(\lambda_1)}.
		\]
		Using Euler's formula, we obtain
		\[
		y_k
		=\frac{\lambda_1^k-\lambda_2^k}
		{2i\,\operatorname{Im}(\lambda_1)}
		=\frac{r^k}{\operatorname{Im}(\lambda_1)}
		\sin(k\phi).
		\]
		By \Cref{sin-lemma}, there exists an integer $k\geq1$ such that
		$\sin(k\phi)\leq0$. Since $r>0$ and
		$\operatorname{Im}(\lambda_1)>0$, it follows that $y_k\leq0$.
	\end{proof}
	
	\begin{lemma}\label{sin-lemma}
		For every $\phi\in\RR$, there exists an integer $k\geq1$ such that
		$\sin(k\phi)\leq0$.
	\end{lemma}
	
	\begin{proof}
		If $\sin\phi\leq0$, take $k=1$. Otherwise, by periodicity, we may
		replace $\phi$ by its representative $\theta\in(0,\pi)$ modulo
		$2\pi$. Set
		$
		k=\left\lceil\frac{\pi}{\theta}\right\rceil.
		$
		Then
		$
		\pi\leq k\theta<\pi+\theta<2\pi,
		$
		so $\sin(k\phi)=\sin(k\theta)\leq0$.
	\end{proof}
	
	
	\begin{example}
		Consider the polynomial
		$
		p(x)=x^2-\frac{19}{10}x+5.
		$
		Its discriminant is
		\[
		\Delta=\left(\frac{19}{10}\right)^2-20
		=-\frac{1639}{100}<0.
		\]
		The associated sequence is given by
		\[
		y_0=0,\qquad y_1=1,\qquad
		y_{k+2}=\frac{19}{50}y_{k+1}-\frac15y_k
		\quad(k\geq0).
		\]
		In particular,
		\[
		y_2=\frac{19}{50},
		\qquad
		y_3=\left(\frac{19}{50}\right)^2-\frac15
		=-\frac{139}{2500}.
		\]
		Thus, $y_3$ is the first nonpositive term after $y_0$.
		We therefore take
		\[
		q(x)=q_2(x)=1+\frac{19}{50}x.
		\]
		We have
		\[
		p(x)q(x)
		=5-5y_3x^2+y_2x^3
		=5+\frac{139}{500}x^2+\frac{19}{50}x^3,
		\]
		which has nonnegative coefficients.
	\end{example}
	
		
\begin{remark}[Historical note]
	Essentially the same construction appears in Meissner's 1911 paper,
	where it is attributed to Hurwitz's 1904 seminar; see
	\cite[Introduction and Footnote~2]{WaymentSteele2023}.
	For the interpretation via truncated reciprocal power series,
	see \cite[Section~3]{KepkaKorbelar2012}.

	In \cite{Polya1928}, Pólya proves the statement and extends it
	to the multivariate case. In particular, he shows that
	$\pm (1+x)^N p(x)$ has nonnegative coefficients for all sufficiently
	large $N$. Explicit bounds on $N$ were later obtained, for example,
	in \cite{Powers2001}.
\end{remark}
	
	\printbibliography
	
\end{document}
